\documentclass[11pt]{article}
\usepackage[a4paper,margin=25mm]{geometry}
\usepackage{amsmath,amssymb,amsthm}
\usepackage[hidelinks]{hyperref}
\newtheorem{theorem}{Theorem}
\title{Ideal Chains Do Not Measure Bratteli Path Depth:\\A Disproof of Conjecture 00000009783}
\author{GitHub: gaochengzhecpu}
\date{}
\begin{document}
\maketitle
\begin{abstract}
The AF algebra $\mathbb C$ has exactly two closed ideals and hence a
unique maximal ideal chain, of length one. Its standard constant
one-dimensional inductive system has a Bratteli diagram with one vertex
at each level and one edge between consecutive levels. This diagram
has paths of every finite length. The claimed equality between ideal
chain length and the supremum of path depths therefore fails. The
counterexample remains valid if chain length counts elements instead
of strict inclusions. Lean verifies the actual complex algebra, its
finite-dimensional approximation and direct-limit property, its closed
ideals, and genuine paths in the associated diagram.
\end{abstract}

\section*{Reading of the source}
Both language versions assert that the length of maximal ideal chains
in an AF algebra equals the supremum of path depths of its Bratteli
diagram. We use the standard meanings: closed two-sided ideals ordered
by inclusion, and depth of a finite directed path measured by its
number of edges. Counting vertices changes path length by one and
does not change whether depths are unbounded. Counting the elements
of an ideal chain instead of its strict inclusions also leaves the
counterexample intact.

An AF algebra is a norm closure of an increasing union of
finite-dimensional $C^*$-subalgebras, equivalently an inductive limit
of finite-dimensional $C^*$-algebras. Finite-dimensional algebras are
included: their constant identity systems are allowed. A Bratteli
diagram has a vertex for each simple matrix summand at each stage,
and the number of edges between successive-stage vertices is the
partial embedding multiplicity. These conventions are stated, for
example, in Blackadar's author-hosted text, sections I.3.2.1--2 and I.3.2.6--8
\cite{blackadar}.

We refute the chain-depth clause. We do not need to choose between
conventions on hereditary subsets or their saturation, nor make any
assertion about the separate dimension-vector stratification clause.
Changing those conventions cannot change the two ideals of the
algebra or the unbounded path lengths of the specified diagram.

\begin{theorem}
There is an AF algebra with a Bratteli diagram whose finite path
depths are unbounded, but whose unique maximal closed-ideal chain
has two elements and one strict inclusion.
\end{theorem}
\begin{proof}
Take $A=\mathbb C$, with its usual norm, multiplication and complex
conjugation. It is a unital complex $C^*$-algebra. Put
\[
 A_n=\mathbb C=M_1(\mathbb C),\qquad
 \varphi_n=\operatorname{id}_{\mathbb C}\quad(n\geq0).
\]
Internally all stages are the entire closed star subalgebra of $A$.
They are one-dimensional, form an increasing sequence, and their
union is exactly $A$. Hence $A$ is AF.

For clarity, this system really has $A$ as its inductive limit.
The canonical maps from all stages into $A$ are identities. For
any $C^*$-algebra $B$, a compatible family of star algebra maps
$f_n:\mathbb C\to B$ satisfies $f_{n+1}=f_n$. It therefore factors
through the unique map $f_0:A\to B$. These maps are continuous:
they are complex-linear maps with finite-dimensional domain.
The same argument also applies to nonunital star algebra maps and
nonunital targets. This is the required direct-limit universal property.

The explicit identification $\mathbb C\to M_1(\mathbb C)$ sends
$z$ to the matrix $[z]$ and respects the star operation. Under this
identification, each connecting map is the identity on $M_1(\mathbb C)$.
It contains exactly one copy of the source matrix block, so its
multiplicity is one. There is consequently one vertex $v_n$ at each
level and exactly one edge $v_n\to v_{n+1}$:
\[
 v_0\longrightarrow v_1\longrightarrow v_2\longrightarrow\cdots.
\]
For every $N\geq0$, the path from $v_0$ through $v_1,\ldots,v_N$
has exactly $N$ edges. Thus the supremum of finite path depths is
$+\infty$.

An ideal $I$ of $\mathbb C$ containing a nonzero element $z$ contains
$z^{-1}z=1$, so it is all of $\mathbb C$. Therefore the only ideals
are $\{0\}$ and $\mathbb C$, and both are closed in the norm topology.
The complete closed-ideal lattice is the two-element chain
\[
 \{0\}\ \subsetneq\ \mathbb C.
\]
It is the unique maximal chain, since every subset can be extended
to this whole two-element chain. Its length is one by strict
inclusions, or two by elements. Neither finite value equals
$+\infty$. This disproves the asserted equality under either
chain-length convention.
\end{proof}

\section*{Formal verification and scope}
The Lean file uses Mathlib's actual complex numbers with the
\texttt{CommCStarAlgebra} instance. Its predicate
\texttt{ApproximatelyFinite} requires a monotone sequence of actual
\texttt{StarSubalgebra} objects, each closed and finite-dimensional,
whose union is dense. The counterexample satisfies it with every
stage equal to the top subalgebra. The actual stage types are
identified with $M_1(\mathbb C)$ by star algebra equivalences.

The connecting maps are actual \texttt{StarAlgHom} identities.
The theorem
\begin{center}\small\texttt{constant\_system\_universal}\end{center}
proves existence and uniqueness of the factorization for every
unital $C^*$-algebra target and compatible unital cocone; continuity
of each map out of a stage is proved separately. A second theorem
proves the universal property, including continuity, for nonunital
$C^*$-algebra targets and nonunital star algebra maps as well.
The internal stage
inclusions and external identity system are checked to agree.
Conjugating the connecting map by the matrix equivalence gives the
actual identity matrix-block embedding. In this one-dimensional
case the multiplicity is the complex dimension of its range,
which is proved to be one. No general formula identifying the
dimension of an arbitrary matrix-algebra image with a partial
embedding multiplicity is asserted.

The diagram is a genuine Mathlib \texttt{Quiver} on $\mathbb N$:
its arrow type from $n$ to $m$ is a finite type with the computed
multiplicity when $m=n+1$, and is empty otherwise. The recursively
constructed \texttt{Quiver.Path} from $0$ to $N$ has length $N$.
The supremum over lengths of all actual paths is taken in
$\mathbb N\cup\{+\infty\}$ and is proved to equal $+\infty$.

\texttt{ClosedIdeal} consists of actual Mathlib ideals with a proof
that their carriers are topologically closed. Since the algebra is
commutative, these are its closed two-sided ideals. Lean proves
that each is the zero or whole ideal, that these are distinct, and
that there are exactly two. It also proves that every maximal chain,
using Mathlib's \texttt{IsMaxChain}, is the whole ideal set.

The maximum strict-inclusion length is encoded by the supremum of
$n$ for which there is a strictly increasing map from
$\{0,\ldots,n\}$ into the closed-ideal order. Lean proves that this
supremum is one, and that the corresponding element-count supremum
is two. These are computations of chains in the actual ideal
order, not stipulated numerical invariants.

The proposition \texttt{OneBlockDepthClause} quantifies over all
commutative complex $C^*$-algebras that are AF and star-algebra
isomorphic to $\mathbb C$. It demands that their ideal-chain length
equal the depth of the constant one-block diagram. Each such
algebra admits the same identity-system presentation. This is a
necessary restriction of the source's claim about all AF algebras;
its negation is the main theorem
\begin{center}\small\texttt{conjecture\_9783\_false}.\end{center}
The alternate element-count convention is separately refuted.

No general AF classification theorem, hereditary-subset
correspondence, or dimension-vector stratification theorem is
claimed as formalized. The exact counterexample, the direct-limit
property, the finite matrix-block identification, the ideals and
the paths are all verified. The Lean file uses no \texttt{sorry},
\texttt{admit}, \texttt{native\_decide}, or custom axiom. Its axiom
reports contain only \texttt{propext}, \texttt{Classical.choice}
and \texttt{Quot.sound}.

\section*{Source and reproduction}
\texttt{SOURCE.md} is an exact copy of the bilingual upstream file
\begin{center}\small\texttt{conjectures/00000009783.md}.\end{center}
The provenance record gives the upstream revision and SHA-256.
The project pins Lean 4.19.0 and Mathlib revision
\begin{center}\small\texttt{c44e0c8ee63ca166450922a373c7409c5d26b00b}.\end{center}
With the pinned dependencies available, run from \texttt{lean/}
\texttt{lake build} and
\begin{center}\small\texttt{lake env lean -DwarningAsError=true Main.lean}.\end{center}
Compile this document with Tectonic or a compatible LaTeX engine.
Verification records distinguish author checks from coordinator
checks. Developed by a delegated AI agent with a self-review;
the coordinating agent reviews it separately. No external
independent review is claimed.

\begin{thebibliography}{1}
\bibitem{blackadar} B. Blackadar, \emph{Classification of
$C^*$-Algebras}, incomplete preliminary edition, June 4, 2025,
sections I.3.2.1--2 (AF algebras) and I.3.2.6--8 (multiplicities and diagrams).
\href{https://bruceblackadar.com/Mathematics/Class.pdf}{Author's public text}.
\end{thebibliography}
\end{document}
